% !TEX root=./main.tex

%%%%%%%%%%%%%%%%%%%%%%%%%%%%%%%% *常用命令 %%%%%%%%%%%%%%%%%%%%%%%%%%%%%%%%
%todo 空格: 1:\  2:~ 3:\quad \qquad 4:\hspace{l}

%todo 分段: \par

%todo 脚注: \footnote{}

%todo 加粗: \textbf{}

%todo 斜体: \textit{}

%todo 行间距: \vspace{0.5\baselineskip} 

%todo 强制行间距: \vspace*{0.5\baselineskip}

%todo \noindent            % 单独取消当前段落的首行缩进

%todo \href{URL}{name}     % 超链接

%todo \setcounter{page}{0} %更改页码

%todo \includepdf[pages=1]{example.pdf}  % 插入pdf第一页

%todo \tableofcontents      % 目录

%todo 引用
%todo \autoref{tab:example}
%todo \ref{tab:example}
%todo \ref{fig:example}

%todo 分栏
%todo \begin{multicols}{n}
%todo n为栏数，代码复制到需要分栏的地方
%todo \end{multicols}

%%%%%%%%%%%%%%%%%%%%% *封面，列表，表格，页眉，代码，图片 %%%%%%%%%%%%%%%%%%%%%

%%%%%%%%%%%%%%% ?封面 %%%%%%%%%%%%%%%

%* 常用 
% \title{标题}
% \author{作者}
% \date{时间}

%* article专属
% \begin{abstract}
% \end{abstract}

%%%%%%%%%%%%%%% ?列表 %%%%%%%%%%%%%%%

%* 这是一个计数的列表
% leftmargin: 控制整体左边距
% itemsep: 控制各item项之间的垂直间距
% topsep: 控制列表与上下文的垂直间距
% label={\rm (\arabic*)} %\roman

% \begin{enumerate}  
% 	\item 第一项
% 		\begin{enumerate}
% 			\item 第一项中的第一项
% 			\item 第一项中的第二项
% 		\end{enumerate}
% 	\item 第二项
% \end{enumerate}

%* 这是一个不计数的列表
% \begin{itemize}  %[label={$\bullet$}]
% 	\item 第一项
% 	\begin{itemize}
% 		\item 第一项中的第一项
% 		\item 第一项中的第二项
% 	\end{itemize}
% 	\item 第二项
% \end{itemize}

%%%%%%%%%%%%%%% ?表格 %%%%%%%%%%%%%%%

% \begin{tabular}{@{}m{5cm}@{}m{5cm}@{}}
%     \raggedright 第一段内容 & 第二段内容 \\[2em]
%     顶对齐测试 \vfill & 底对齐测试 \vfill \\
% \end{tabular}

% 带边框的表格
% m表示该列内容宽度也是 5 厘米，但内容在垂直方向上会居中对齐
% p表示该列内容的宽度是 5 厘米，内容会自动换行
% | 表示列之间有竖线 \hline横线 

% \begin{tabular}{|p{5cm}|m{5cm}|}   
%     \hline
%     上对齐 & \parbox[t]{5cm}{这是顶对齐的内容，\\支持多行文本展示。} \\ \hline
%     居中对齐 & \parbox[c]{5cm}{这是居中对齐的内容，\\支持多行文本展示。} \\ \hline
%     下对齐 & \parbox[b]{5cm}{这是底对齐的内容，\\支持多行文本展示。} \\ \hline
% \end{tabular}

%%%%%%%%%%%%%%% ?代码 %%%%%%%%%%%%%%%

% \lstinputlisting[
%     language    =   Python,
%     caption     =   {\bf code.py},
%     label       =   {code.py}
% ]{./codes/code.py}

%%%%%%%%%%%%%%% ?图片 %%%%%%%%%%%%%%%
%%%%% *单独一张 %%%%%
% \begin{figure}[htbp]  %? h:当前 t:顶部 b:底部 p:单独一页 H:固定
%     \flushleft          %? 左对齐
%     \centering          %? 居中
%     \zihao{-5}
%     \songti
%     \includegraphics[width=0.5\textwidth]{相对路径}  %宽度为纸张的一半
%     \caption{标题}
%     \label{fig:标签}
% \end{figure}

%%%%% *并排几张 %%%%%

%? 也可以调整高度 [height=0.2]

% \begin{figure}[ht]
%     \centering
%     \begin{minipage}{0.4\textwidth}
%         \centering
%         \includegraphics[width=\textwidth]{image1.jpg}
%         \caption{第一张图片}
%         \label{fig:img1}
%     \end{minipage}
%     \hfill
%     \begin{minipage}{0.4\textwidth}
%         \centering
%         \includegraphics[width=\textwidth]{image2.jpg}
%         \caption{第二张图片}
%         \label{fig:img2}
%     \end{minipage}
% \end{figure}

%%%%% *环绕文字 %%%%%
% \begin{wrapfigure}{r}{0.4\textwidth} % r 表示图片在右侧，0.4\textwidth 是图片宽度
%     \centering
%     \includegraphics[width=0.4\textwidth]{example.jpg}
%     \caption{这是环绕图片的标题}
%     \label{fig:example}
% \end{wrapfigure}

%%%%%%%%%%%%%%% ?三线表 %%%%%%%%%%%%%%%

% \begin{table}[ht]
%     \centering
%     \caption{示例三线表}
%     \label{tab:example}
 
%     \begin{tabular}{@{} l c r @{}} %* @{} 去掉左右多余的 \tabcolsep
%       \toprule
%       左对齐 & 居中 & 右对齐 \\
%       \midrule
%       第一行 & 数据 A & 123 \\
%       第二行 & 数据 B & 456 \\
%       第三行 & 数据 C & 789 \\
%       \bottomrule
%     \end{tabular}
% \end{table}

%%%%%%%%%%%%%%% ?公式 %%%%%%%%%%%%%%%
%* 行内公式
% $ \theta $

%* 行间公式
% \begin{equation}
% E=mc^2
% \label{eq:energy}
% \end{equation}

%* 单行超长公式
% \begin{multline}
% a+b+c+d+e+f+g+h+i+j+k+l+m+n+o+p+q+r+s+t+u+v+w+x+y+z \\
% = 1+2+3+4+5+6+7+8+9+10+11+12+13+14+15+16+17+18+19+20
% \end{multline}

%* 多行公式
% \begin{align}
% P & = UI \\
% & = I^2R
% \end{align}

%* 矩阵
%* 可换成equation来添加序号
% \[    
% \mathbf{X} = \left(
%     \begin{array}{cccc}
%     x_{11} & x_{12} & \ldots & x_{1n}\\
%     x_{21} & x_{22} & \ldots & x_{2n}\\
%     \vdots & \vdots & \ddots & \vdots\\
%     x_{n1} & x_{n2} & \ldots & x_{nn}\\
%     \end{array} \right)
% \]

%* 分段函数
% \[
% f(x) =
%     \begin{cases}
%         0 &  x \text{为无理数} ,\\
%         1 &  x \text{为有理数} .
%     \end{cases}
% \]

%* 公式内加粗
% \bm{}

%* 宏包手册
% http://texdoc.net/texmf-dist/doc/latex/amsmath/amsldoc.pdf
% http://mirrors.ctan.org/info/symbols/comprehensive/symbols-a4.pdf

%%%%%%%%%%%%%%% ?定义 %%%%%%%%%%%%%%%
% \begin{definition}
%     定义的内容
%     \label{def:nosense}
% \end{definition}

% definition   def  定义
% theorem      thm  定理
% lemma        lem  引理
% corollary    cor  推论
% assumption   asu  假设
% conjecture   con  猜想
% axiom        axi  公理
% principle    pri  定律
% problem      pro  问题
% example      exa  例 
% solution     sol  解 

%* 用ref 别用autoref